\documentclass[12pt]{article}

\usepackage[utf8]{inputenc}
\usepackage{amsmath, amssymb, amsthm}
\usepackage{geometry}
\geometry{a4paper, margin=2.5cm}

\title{Spațiu Măsurabil}
\author{}
\date{}

\theoremstyle{definition}
\newtheorem{definition}{Definiție}

\begin{document}

\maketitle

\begin{definition}
Fie $\Omega$ un spațiu și $\mathcal{F}$ o sigma-algebră de submulțimi ale lui $\Omega$. 
Perechea $(\Omega, \mathcal{F})$ se numește \emph{spațiu măsurabil}.
\end{definition}

\noindent
Elementele lui $\mathcal{F}$ se numesc \emph{mulțimi măsurabile}.

\end{document}
